\documentclass[12pt,a4paper]{article}

% ------------------------------------------------
% PACKAGES
% ------------------------------------------------
\usepackage[margin=1in]{geometry}
\usepackage{setspace}
\usepackage{titlesec}
\usepackage{enumitem}
\usepackage{booktabs}
\usepackage{longtable}
\usepackage{array}
\usepackage{hyperref}
\usepackage{fancyhdr}
\usepackage{graphicx}

\hypersetup{
    colorlinks=true,
    linkcolor=black,
    urlcolor=blue
}

\pagestyle{fancy}
\fancyhf{}
\rhead{IBDP Psychology}
\lhead{Paper 1 Guide}
\rfoot{\thepage}

\setstretch{1.15}

% ------------------------------------------------
% TITLE
% ------------------------------------------------
\title{\textbf{IBDP Psychology \\ Paper 1 – Teacher Guide and Assessment Framework}}
\author{Internal Academic Resource}
\date{}

\begin{document}

\maketitle
\thispagestyle{empty}

\vspace{1cm}

\section*{Purpose of This Document}
This document provides a structured overview of Paper 1 in IBDP Psychology. 
It is intended for instructional planning, student preparation, and internal academic review. 
This guide paraphrases official structures and does not reproduce proprietary assessment material.

\newpage

% ------------------------------------------------
\section{Paper 1 Overview}

\begin{itemize}
    \item \textbf{Total Marks:} 35
    \item \textbf{Duration:} 1 hour 30 minutes
    \item \textbf{Weighting:} 
    \begin{itemize}
        \item Standard Level (SL): 35\%
        \item Higher Level (HL): 25\%
    \end{itemize}
\end{itemize}

\subsection*{Focus of Assessment}
Paper 1 assesses integration of:
\begin{itemize}
    \item Core psychological approaches (biological, cognitive, sociocultural)
    \item Research methodology
    \item Conceptual understanding
    \item Application within specified contexts
\end{itemize}

\subsection*{Key Concepts}
Students are expected to engage with six central concepts:
\begin{itemize}
    \item Bias
    \item Causality
    \item Change
    \item Measurement
    \item Perspective
    \item Responsibility
\end{itemize}

\subsection*{Contexts}
Questions may be framed within:
\begin{itemize}
    \item Health and well-being
    \item Human development
    \item Human relationships
    \item Learning and cognition
\end{itemize}

\newpage

% ------------------------------------------------
\section{Section A: Short Answer Questions (4 + 4 marks)}

\subsection*{Structure}
\begin{itemize}
    \item Two compulsory questions
    \item Each worth 4 marks
    \item Drawn from two of the three approaches
\end{itemize}

\subsection*{Skills Assessed}
\begin{itemize}
    \item Knowledge and understanding (AO1)
    \item Application and analysis (AO2)
\end{itemize}

\subsection*{What Students Must Demonstrate}
\begin{itemize}
    \item Clear explanation of a theory or concept
    \item Accurate psychological terminology
    \item A relevant and explained example
\end{itemize}

\subsection*{Performance Expectations}

\begin{longtable}{p{1.5cm} p{12cm}}
\toprule
\textbf{Marks} & \textbf{General Performance Description} \\
\midrule
0 & Response does not meet minimum standard. \\
1--2 & Basic or limited explanation; example may be present but insufficiently developed. \\
3--4 & Clear, accurate explanation; example is relevant and well linked to the concept. \\
\bottomrule
\end{longtable}

\subsection*{Teacher Strategy}
\begin{itemize}
    \item Train students to define first, then explain, then exemplify.
    \item Avoid unnecessary research study detail unless directly relevant.
    \item Encourage concise structure (approx. 8–10 sentences).
\end{itemize}

\newpage

% ------------------------------------------------
\section{Section B: Application Questions (6 + 6 marks)}

\subsection*{Structure}
\begin{itemize}
    \item Two compulsory questions
    \item Each worth 6 marks
    \item Based on unseen scenarios
    \item Each question linked to a specific context
\end{itemize}

\subsection*{Skills Assessed}
\begin{itemize}
    \item Application of knowledge (AO2)
    \item Analytical reasoning
\end{itemize}

\subsection*{What Students Must Demonstrate}
\begin{itemize}
    \item Accurate understanding of relevant theory
    \item Direct application to scenario details
    \item Logical explanation of how theory explains or addresses situation
\end{itemize}

\subsection*{Performance Expectations}

\begin{longtable}{p{1.5cm} p{12cm}}
\toprule
\textbf{Marks} & \textbf{General Performance Description} \\
\midrule
0 & Response does not meet minimum standard. \\
1--2 & Limited understanding; minimal or weak application. \\
3--4 & Mostly accurate knowledge; partial and relevant application. \\
5--6 & Accurate and detailed knowledge; strong, well-developed application. \\
\bottomrule
\end{longtable}

\subsection*{Teacher Strategy}
\begin{itemize}
    \item Train students to reference specific elements of the scenario.
    \item Use the structure: Identify theory → Explain theory → Apply explicitly to scenario.
    \item Emphasize “link-back” sentences.
\end{itemize}

\newpage

% ------------------------------------------------
\section{Section C: Extended Response (15 marks)}

\subsection*{Structure}
\begin{itemize}
    \item Two concept-based essay questions
    \item Students answer one
    \item Each question linked to a different context
\end{itemize}

\subsection*{Skills Assessed}
\begin{itemize}
    \item Knowledge and understanding (AO1)
    \item Synthesis and evaluation (AO3)
    \item Conceptual integration
\end{itemize}

\subsection*{Expectations for High-Level Responses}
\begin{itemize}
    \item Clear focus on the command term (e.g., evaluate, discuss)
    \item Integration of at least one key concept
    \item Critical analysis (strengths, limitations, implications)
    \item Logical structure and coherent argument
    \item Clear and reasoned conclusion
\end{itemize}

\subsection*{Performance Bands (Generalized)}

\begin{longtable}{p{1.5cm} p{12cm}}
\toprule
\textbf{Marks} & \textbf{General Performance Description} \\
\midrule
1--3 & Limited understanding; largely descriptive; weak terminology. \\
4--6 & Some understanding; limited analysis; simplistic conclusion. \\
7--9 & Adequate understanding; developing analysis; partially sustained argument. \\
10--12 & Clear explanation; relevant analysis; consistent argument and conclusion. \\
13--15 & Fully developed explanation; strong critical analysis; precise terminology; reasoned and coherent conclusion. \\
\bottomrule
\end{longtable}

\subsection*{Recommended Essay Structure}
\begin{enumerate}
    \item Introduction (define key terms + outline argument)
    \item Theoretical explanation
    \item Supporting analysis
    \item Critical evaluation
    \item Concept integration
    \item Conclusion
\end{enumerate}

\newpage

% ------------------------------------------------
\section{Exam Preparation Framework}

\subsection*{Section A Practice Routine}
\begin{itemize}
    \item Weekly 4-mark drills
    \item Focus on concise explanation
    \item Peer marking using band descriptors
\end{itemize}

\subsection*{Section B Practice Routine}
\begin{itemize}
    \item Scenario-based workshops
    \item Structured application templates
    \item Timed practice (12 minutes per question)
\end{itemize}

\subsection*{Section C Practice Routine}
\begin{itemize}
    \item Concept-mapping exercises
    \item Timed essay simulations
    \item Focused evaluation paragraph drills
\end{itemize}

% ------------------------------------------------
\section*{Academic Integrity Notice}

This document is an internally developed instructional guide. 
It summarizes and paraphrases publicly available structural information 
about IBDP Psychology Paper 1 for educational planning purposes. 
It does not reproduce proprietary examination content.

\vfill

\begin{center}
\textit{End of Document}
\end{center}

\end{document}